\section{Innledning}
\label{sec:krav-innledning}

Kravspesifikasjonen ble utarbeidet av HDO ved prosjektstart og 
dannet grunnlaget for all videre utvikling. Kravene ble mottatt 
fra Magne Henriksen og er gjengitt i sin helhet i 
vedlegg~\ref{app:kravspec}~\cite{hdo_kravspec}.

Krav-ID-ene er beholdt fra HDOs opprinnelige kravspesifikasjon for å 
sikre sporbarhet mellom krav, design, implementasjon og testing. Der kravene 
var åpne eller backend-avhengige, beskriver dette kapittelet gruppens tolkning 
og avgrensning. Selve løsningen og graden av kravoppfyllelse vurderes senere 
i design-, implementasjons-, test- og diskusjonskapitlene.

HDOs opprinnelige kravspesifikasjon inneholdt ikke en eksplisitt prioritering av kravene. I prosjektet ble funksjonelle kjerneoperasjoner, sikkerhet og støtte for begge videomotorer behandlet som må-ha-krav, mens observability og dokumentasjon ble behandlet som viktige støttekrav for drift, testing og overlevering.

Kravene er organisert i følgende kategorier:
\begin{itemize}
    \item Funksjonelle krav (F1.1--F1.5)
    \item Sikkerhetskrav (S1.1--S1.2)
    \item Infrastrukturkrav (I1.1--I1.2)
    \item Observability-krav (O1.1--O1.3)
    \item Dokumentasjonskrav (D1.1)
\end{itemize}

\section{Oversikt over krav og evaluering}
\label{sec:krav-oversikt}

\begin{table}[H]
\centering
\begin{threeparttable}
\caption{Oversikt over krav, prioritet og evalueringsgrunnlag.}
\label{tab:krav-oversikt}
\begin{tabular}{l l l l}
\toprule
\textbf{ID} & \textbf{Krav} & \textbf{Prioritet} & \textbf{Evalueres gjennom} \\
\midrule
F1.1 & Opprette videorom & Må ha & Integrasjonstest, E2E-test \\
F1.2 & Slette videorom & Må ha & Integrasjonstest, E2E-test \\
F1.3 & Hente rominformasjon & Må ha & Integrasjonstest, manuell test \\
F1.4 & Starte videostrøm & Backend-avhengig & Provider-test, E2E-test \\
F1.5 & Stoppe videostrøm & Backend-avhengig & Provider-test, E2E-test \\
S1.1 & API-autentisering & Må ha & Sikkerhetstest, integrasjonstest \\
S1.2 & Romspesifikk tilgang & Må ha & E2E-test, dokumentasjonskontroll \\
I1.1 & Containerisering & Må ha & Docker-test, deployment-test \\
I1.2 & Stateless og konfig & Må ha & Arkitekturgjennomgang, container-test \\
O1.1 & Logging & Viktig & Logggjennomgang, manuell test \\
O1.2 & Helsesjekker & Viktig & API-test, E2E-test \\
O1.3 & Prometheus-metrikker & Viktig & Metrics-test, Grafana-verifisering \\
D1.1 & Dokumentasjon & Viktig & Dokumentasjonsgjennomgang \\
\bottomrule
\end{tabular}
\begin{tablenotes}
\footnotesize
\item F1.6 var ikke spesifisert av HDO og er utelatt.
\end{tablenotes}
\end{threeparttable}
\end{table}